\subsection{Biyolojik Temeller}

Walrus Optimizer (WO), morsaların yaşam biçimlerinden ve grup
davranışlarından esinlenen sürü zekâsı tabanlı bir meta-sezgisel
algoritmadır. Han ve arkadaşları, özellikle göç etme, üreme, birlikte
dinlenme, besin arama ve tehlikeden kaçınma davranışlarını algoritmanın
kavramsal temelini oluşturmak için kullanmıştır.

Morsalar, çoğunlukla Arktik çevresindeki soğuk denizlerde yaşayan ve
gruplar hâlinde hareket eden memelilerdir. Karada ve buz üzerinde
dinlenebilir, su altında besin arayabilirler. Üst dudaklarındaki hassas
bıyıklar çevreyi algılamalarına yardımcı olur; uzun köpek dişleri ise
savunma, deniz tabanındaki besinleri ortaya çıkarma ve buz üzerine
tırmanma gibi işlevler görür. Bu özellikler, WO'da kullanılan
mekanizmaların biyolojik arka planını açıklar; ancak algoritmadaki
matematiksel işlemler bu davranışların doğrudan biyolojik kopyaları
değil, optimizasyon amacıyla yapılmış soyutlamalardır.

\paragraph{Sosyal ve çevresel davranışlar.}
WO'nun dayandığı biyolojik gözlemler aşağıdaki şekilde özetlenebilir:

\begin{itemize}
    \item \textbf{Üreme döneminde alan seçimi:} Üreme döneminde erkek
    morsalar sahilde bölgeler edinir. Makalede, daha güçlü erkeklerin
    daha elverişli bölgeleri elde edebildiği ve bölgenin büyüklüğünün
    eşleştikleri dişi sayısıyla ilişkili olduğu belirtilir.

    \item \textbf{Besin arama ve bilgi paylaşımı:} Morsalar deniz
    tabanında besin ararken çevrelerinden yararlanır. Makale, onların
    ses yoluyla yön bulma ve besin bilgilerini diğer bireylerle paylaşma
    davranışlarına dikkat çeker.

    \item \textbf{Toplu savunma ve yardımlaşma:} Katil balina gibi
    tehditlerle karşılaştıklarında morsalar grup hâlinde savunma
    davranışı gösterebilir ve yaralanan bireylere yardım edebilir.
    Bu gözlem, tehlikeye karşı sürü üyelerinin birbirleriyle
    etkileşimini vurgular.

    \item \textbf{Göç ve dinlenme:} Morsalar çevresel koşullara ve
    yaşam döngülerine bağlı olarak hareket eder; karada veya buz
    üzerinde topluca dinlenebilir, suya girerek besin arayabilir.
\end{itemize}

\paragraph{Biyolojik davranışların optimizasyonla ilişkisi.}
WO, bu gözlemleri arama sürecinde farklı işlevleri olan mekanizmalar
olarak ele alır. Göç davranışı, arama ajanlarının çözüm uzayında yeni
bölgelere yönelmesini; besin arama, aday çözümlerin değerlendirilmesini
ve umut vadeden bölgelerde daha ayrıntılı arama yapılmasını temsil eder.
Üreme davranışı, başarılı çözümlerden yararlanarak yeni aday çözümler
oluşturma fikrine karşılık gelir. Dinlenme ve toplu davranışlar sürü
içindeki etkileşimleri, tehlikeden kaçınma ise arama sürecinde konum
değiştirerek yerel çözümlere aşırı bağlı kalma riskini azaltma fikrini
destekler.

Bu biyolojik analojiler, WO'da arama uzayını keşfetme
(\textit{exploration}) ile mevcut iyi çözümlerin çevresinde yoğunlaşma
(\textit{exploitation}) arasında denge kurmak amacıyla kullanılır.
Dolayısıyla morsaların davranışları algoritmanın kavramsal
esin kaynağını oluştururken, gerçek optimizasyon işlemleri makalenin
sonraki bölümünde tanımlanan matematiksel güncelleme kurallarıyla
gerçekleştirilir.